% !PRINTED 300
\weberpage{300}

Es seien wieder $\alpha$ und $\beta > \alpha$ zwei reelle Zahlen, die das Intervall $(\alpha, \beta)$ bestimmen, in dem die Anzahl der reellen Wurzeln von $f(x)$ gezählt werden sollen.

Wir wollen zunächst annehmen, dass in dem Intervall $(\alpha, \beta)$ nicht zwei Glieder der Reihe (1) zugleich verschwinden, und dass für $x = \alpha$, $x = \beta$ kein Glied der Reihe verschwindet.

Lassen wir nun $x$ von $\alpha$ bis $\beta$ stetig wachsen, so kann in der Vorzeichenfolge der Reihe (1) nur dann eine Aenderung eintreten, wenn eine der Functionen durch Null geht. Wenn $f^{(\nu)}(x)$ für $x = \xi$ durch Null geht, so geht $f^{(\nu)}(x)$, je nachdem $f^{(\nu+1)}(x)$ positiv oder negativ ist, von negativen zu positiven oder von positiven zu negativen Werthen über, und es geht also zwischen $f^{(\nu)}$ und $f^{(\nu+1)}$ beim Durchgang durch $\xi$ ein Zeichenwechsel verloren. Dies gilt auch, wenn $f^{(\nu)}$ die Function $f(x)$ selbst ist. Wenn aber $f^{(\nu)}$ eine der Derivirten ist, so geht ihr eine Function $f^{(\nu-1)}$ voran, und da $f^{(\nu-1)}$ nach Voraussetzung für $x = \xi$ nicht Null ist, und $f^{(\nu)}$ beim Durchgang durch $\xi$ sein Zeichen wechselt, so findet zwischen $f^{(\nu-1)}$ und $f^{(\nu)}$ beim Durchgang durch $\xi$ entweder ein Verlust oder ein Gewinn von einem Zeichenwechsel statt. Wenn also ein inneres Glied der Reihe (1) durch Null geht, so bleibt die Anzahl der Zeichenwechsel ungeändert, oder es gehen zwei Zeichenwechsel verloren.

Geht aber $f(x)$ selbst durch Null, so geht ein Zeichenwechsel verloren.

Daraus folgt das Theorem:

I. \gesperrt{Die Anzahl der zwischen $\alpha$ und $\beta$ gelegenen Wurzeln von $f(x)$ ist höchstens so gross, wie die Zahl der zwischen $\alpha$ und $\beta$ verlorenen Zeichenwechsel, und wenn sie kleiner ist, so ist der Unterschied eine gerade Zahl.}

Mit Benutzung einer Formel können wir auch sagen:

Ist $V(x)$ die Anzahl der Zeichenwechsel (Variationen), die die Reihe (1) für irgend einen Werth $x$ darbietet, so ist die Anzahl der zwischen $\alpha$ und $\beta$ gelegenen Wurzeln von $f(x)$

\weq{2}{V(\alpha) - V(\beta) - 2h,}

worin $h$ eine nicht negative ganze Zahl ist.

Der Beweis des Satzes I. bedarf noch einer Ergänzung für den Fall, dass in der Reihe (1) mehrere auf einander folgende Glieder zugleich verschwinden.
